\subsection{15.76(b): automorphisms of categories of free soluble groups}
\label{cand:15.76}
Let \(\mathcal S_d\) be the full variety of groups of derived length
at most \(d\ge1\). Every automorphism of its category of finite-rank
free groups is naturally isomorphic to the identity.
This includes the full metabelian variety. It does not assert the same
for every proper subvariety of soluble groups: the proper metabelian
variety studied by Fernandes--Tsurkov provides a counterexample
\cite{fernandes-tsurkov}.
The reduction of category automorphisms to verbal operations is an
established method, used in the nilpotent predecessor
\cite{tsurkov-nilpotent}. The additional step here is the affine
coefficient induction for the full soluble varieties.

We first recall the affine coefficient model. For an absolutely free
group \(F=F(x_1,\ldots,x_n)\), a normal subgroup \(N\), and \(Q=F/N\),
map \(x_i\) to \((\bar x_i,e_i)\) in
\[
 Q\ltimes\Z[Q]^n,\qquad(q,U)(r,V)=(qr,U+qV).
\]
Its kernel is \([N,N]\). Indeed the covering graph of the rose
corresponding to \(N\) has vertices \(Q\) and edges
\(q\to q\bar x_i\). The two affine coordinates of a word are its
endpoint and signed integral edge chain. A loop has zero chain
exactly when its class in the first homology of this graph is zero,
which is its image in \(N/[N,N]\).
Finite support makes the argument valid for infinite covering graphs.

For a word \(w(x,y)\), write its affine value as
\[
 w((a,U),(b,V))=(w(a,b),f(a,b)U+g(a,b)V).
\]
The boundary of its edge chain gives the ordered group-ring identity
\[
 f(a,b)(a-1)+g(a,b)(b-1)=w(a,b)-1.                 \tag{b}
\]
We prove by induction on \(d\) that a normalized associative word in
\(\mathcal S_d\), meaning
\[
 w(x,1)=x,\quad w(1,y)=y,\quad
 w(w(x,y),z)=w(x,w(y,z)),
\]
must equal \(xy\) or \(yx\) as an identity.
For \(d=1\), exponent sums suffice.
For \(d\ge2\), induction and, if necessary, reversal of the two
arguments allow us to assume \(w(a,b)=ab\) in \(\mathcal S_{d-1}\).
Let \(Q\) be free in that variety on \(a,b,c\).
The affine group \(Q\ltimes\Z[Q]^3\) belongs to \(\mathcal S_d\).
Associativity, evaluated at its three module basis vectors, gives
\[
 f(ab,c)f(a,b)=f(a,bc),\qquad
 g(ab,c)=g(a,bc)g(b,c).
\]
Set \(h(t)=f(1,t)\in\Z[t,t^{-1}]\).
The first equation gives \(f(b,c)h(b)=h(bc)\);
putting \(c=b^{-1}\) shows that \(h\) is a Laurent unit with
\(h(1)=1\). Hence \(h=t^k\) for an integer \(k\), and
\[
 f(a,b)=(ab)^ka^{-k}.
\]
Only the units of the ordinary cyclic Laurent ring are used:
highest and lowest powers in a product show that a unit is
\(\pm t^k\). Likewise, using \(q(t)=g(t,1)\), the second equation gives
\[
 g(a,b)=(ab)^lb^{-l}
\]
for an integer \(l\).
Map (b) to the free abelian quotient, with variables \(X,Y\).
It becomes
\[
 Y^k(X-1)+X^l(Y-1)=XY-1.
\]
Differentiate in \(X\) and set \(X=1\), obtaining
\(Y^k=l+(1-l)Y\).
Coefficient comparison gives exactly \((k,l)=(0,1)\) or \((1,0)\).

In the first case \(f=1,g=a\), the affine coefficients of \(xy\).
Apply the kernel model to \(F(x,y)\) with \(N=F^{(d-1)}\):
equality of affine values proves \(w=xy\) modulo \(F^{(d)}\).
In the second case \(f=aba^{-1},g=1\); equation (b) forces
\(aba^{-1}=b\) in \(Q\).
For \(d\ge3\), \(Q\) maps onto the nonabelian metabelian group \(S_3\),
so this is impossible. For \(d=2\), \(Q\) is abelian and these are
exactly the affine coefficients and scalar coordinate of \(yx\).
The same kernel model gives \(w=yx\).
Undoing any initial reversal proves the induction.

We finish the categorical reduction explicitly.
Rank one is recognized by its infinite commutative endomorphism
monoid: rank zero has a one-element monoid, and ranks at least two
have noncommuting basis projection and interchange endomorphisms.
Coproducts of rank-one objects then recognize every finite rank.
Conjugate the given functor by chosen object isomorphisms to fix
objects, and by basis automorphisms to fix the maps
\(\iota_i:F_1\to F_n\) selecting the basis. Their images form a
coproduct diagram and hence a free basis, justifying the second step.
At \(F_1\) this normalization is the identity.

For \(a\in F_n\), let \(f_a:F_1\to F_n\) send the generator to \(a\),
and set \(s_n(a)=\Psi(f_a)(x)\).
These are bijections fixing the basis and identity.
Functoriality gives, for every homomorphism \(u:F_n\to F_m\),
\[
 \Psi(u)=s_mus_n^{-1}.
\]
With \(w=s_2(xy)\), it also gives
\(s_n(ab)=w(s_n(a),s_n(b))\).
Thus \(w\) transports a group law and is normalized associative on
all free objects, hence throughout \(\mathcal S_d\).
The word result applies.
If \(w=xy\), every \(s_n\) is a basis-fixing homomorphism and is
the identity. If \(w=yx\), let \(\eta_n\) invert each basis generator
as an ordinary automorphism, and let \(j_n\) be the inversion set map.
Then \(s_n=j_n\eta_n\).
Every group homomorphism commutes with the inversion maps, so
\(\Psi(u)=\eta_mu\eta_n^{-1}\), an inner functor.
Undoing the earlier object and basis conjugations supplies a natural
isomorphism for the original functor on every object and morphism.
The argument works equally with all named free objects or a skeleton;
rank zero contributes its unique maps.
The affine kernel model is the classical Magnus method, proved here
from graph homology. No general group-ring unit theorem or classification
of automorphisms of free soluble groups is imported.
See \artifact{research/15.76-proof.md}.
